% ===== §12 The matter of the field: the everyday sensory-habitual fabric =====
\section{The Matter of the Field, and the Conditions of a Just Eudaimonia}
\label{sec:fieldmatter}

\noindent
The field has been specified structurally, as a curvature that guides derivations rather than a structure that is constructed (\xref{sec:field-not-grammar}); it has not been specified materially. This section supplies the matter of the field, and thereby brings the practice to the ground of experience. It identifies the matter of the field with the everyday sensory and habitual fabric of a shared life, names that fabric in the terms of the traditions that have addressed it, shows that the fabric is transversal across the three registers, and specifies the two conditions a field must satisfy to be good. It is the concluding section of the paper's account of practice, and it marks the point at which the account passes to the more empirical register that a subsequent paper will develop.

\subsection{The matter of the field}
\label{sec:fm-matter}

The question the structural account leaves open is the question of what the field is made of. A curvature is a curvature of something; a field that guides the derivations of a shared life is a field constituted by some material. The material is the everyday sensory and habitual fabric of the shared life: the tastes of the food two persons eat together, the sounds of their common rooms, the scents of their shared space, the textures of their contact, and the habits, customs, memories, and household manners that order these. The field is not an abstract curvature laid over a relation; it is this fabric, jointly sustained, within which the relation's derivations are taken.

This fabric is, in the framework of this paper, the site at which the three registers are coupled. Various traditions of practical and contemplative thought have attended to the sensory and habitual fabric of a dwelling and a shared life under their own descriptions; the present account does not depend on any of them, and notes only that the matter it identifies is not a novel object but one long recognised, under other names, as the medium in which a form of life is borne.

\subsection{The fabric is transversal across the registers}
\label{sec:fm-transversal}

The fabric is the matter of the field because it is precisely where the coupling of the three registers, on which value is emergent (\xref{sec:emergence}), is materially sustained. A meal two persons have eaten together for years is not a fact of one register. It has a real aspect, the taste and the bodily satisfaction and the inarticulable ease that no description exhausts; a symbolic aspect, the meal as a sign that is theirs, the mark of a shared life; and an imaginary aspect, the idealised form of the life the meal figures. The meal, and the habit of it, is the coupling of the three registers materially deposited in a recurrent sensory event. So too a shared waking, a way of keeping a room, a manner of touch, a household custom: each is a site at which the three registers are coupled in a sensory-habitual recurrence.

\begin{claim}[The matter of the field is the coupling of the registers deposited in the sensory-habitual fabric]
The matter of the field is the everyday sensory-habitual fabric of a shared life, its forms, sounds, scents, tastes, textures, and the habits, rules, and memories that order them, in which the coupling of the three registers, on which relational value is emergent, is materially sustained. The fabric is transversal across the registers: each of its elements, a shared meal, a common habit, a household custom, has a real, a symbolic, and an imaginary aspect, and is the site at which the three are coupled in a recurrent sensory event. To cultivate the field is to tend this fabric such that the coupling is sustained; and the obstruction or severance of the coupling, in the fabric, is the material form of the settling that extinguishes value (\xref{sec:nameless-corrected}).
\end{claim}

\noindent
This is the material content of the practice. The customs and habits that compose a shared life, the household manners, the recurrent sensory events, are the medium through which a relation's grammar is jointly sustained and rewritten, and they are, in the terms of the practical traditions, the dispositional fabric, the \emph{habitus}, in which a form of life is borne \citep{bourdieu1990}. The cultivation of the field is the tending of this fabric, and the body of practical knowledge concerning how the fabric is tended, in the formation of a household manner and in the small recurrent practices of a shared life, is the empirical content of the practice that the structural account abstracts from.

\subsection{The two conditions of a good field}
\label{sec:fm-conditions}

A good field must satisfy two conditions, and they are the two orthogonal criteria of \xref{sec:criterion}, now stated as conditions on the fabric. The first is eudaimonic: the field must be one in which each party's own dynamics can unfold, in which the sensory-habitual fabric guides the derivations of each toward the development of that party's own powers rather than confining them to a foreseen form. The second is just: the field must produce a good cycle of value, one in which no party is reduced to the fuel of the fabric, present only as the supplier of a sensory-habitual order from which it is never positively reconfigured in return.

The relation between the two conditions requires precision, for the eudaimonic condition is not sufficient and is not, by itself, the criterion. A field may permit the unfolding of one party's dynamics at the cost of another's, as when the sensory-habitual fabric of a household is sustained by the uncompensated labour of one party, characteristically the party who supplies the meals, the order, the atmosphere, the care, and whose own dynamics is thereby confined to the supplying. Such a field produces the satisfaction of one party and the depletion of another, and it is not good, the satisfaction of the one being purchased by the symbolic and material exploitation of the other (\xref{sec:exploitation}). The eudaimonic condition is therefore not self-standing; it is constrained by the just.

\begin{claim}[Eudaimonia is constituted, not merely constrained, by justice]
The just is not an external constraint upon the eudaimonic but internal to it. Eudaimonia is the unfolding of a party's own dynamics; and a field that depletes a party, that reduces it to the fuel of the fabric and confines its dynamics to the supplying, does not produce that party's eudaimonia but forecloses it. A field that produces the satisfaction of one party by the depletion of another has therefore not produced eudaimonia and added an injustice; it has produced no eudaimonia at all, but the enjoyment of one and the foreclosure of the other. Eudaimonia, being the unfolding of each party's dynamics, is irreducibly common, and a satisfaction that is not common, that is purchased by the foreclosure of another's unfolding, is not eudaimonia but enjoyment at another's cost. The just field is not the eudaimonic field with a constraint added; the unjust field is one that has failed to be eudaimonic at all. Justice is the condition under which there is eudaimonia, and not a limit upon it.
\end{claim}

\subsection{Happiness is not the criterion}
\label{sec:fm-nottruth}

A final qualification completes the account and connects it to the ethics. The good field is not to be constructed as a procedure for the guarantee of happiness, and happiness is not the paper's final criterion. The eudaimonia a good field produces is, like the good itself (\xref{sec:keystone}), unguaranteeable and processual: it is not a state that a correctly constructed field secures, but an unfolding that a tended field does not obstruct. To treat the field as a technique for the production of happiness, to construct the sensory-habitual fabric as a machine for a guaranteed satisfaction, is the error of possession in the material register, the promise to engineer a guaranteed fortune from the arrangement of the fabric, and it is, like every construction of structure, the occupation of the place of the Other (\xref{sec:exploitation}). The field is cultivated, not constructed; the happiness is undergone, not secured; and the fabric is tended such that a good cycle of value may generate, not engineered such that a specified satisfaction must result.

This is the conclusion of the paper's account of practice. The practice is the cultivation of a field; the field is the everyday sensory-habitual fabric of a shared life in which the three registers are coupled; the good field satisfies two conditions, the eudaimonic and the just, of which the just is the condition of the eudaimonic rather than a limit upon it; and the cultivation is a tending and not a construction, a guidance of the fabric's derivations toward a good cycle that it may not possess and must not guarantee. The empirical development of this account, the formation of household manners and habits, the practical cultivation of the sensory-habitual fabric, and the theory and practice of a just eudaimonia in the more empirical register, is the matter of the subsequent paper, and the present paper, having grounded the practice in the matter of the field, concludes its account here.
